\documentclass{article}
\usepackage{icmlko}
\usepackage{graphicx}
\usepackage{booktabs}
\usepackage{amsmath}

\title{한 칸을 통째로 빼고도 이겼다 \\
  {\large 효과의 \textbf{절반}(49.05\%)을 나르는 첫 간격 칸을 없애도 순열이 진다 ---
   그런데 이번에는 천장을 \emph{LOO} 로 쟀고, 그래도 위에 있었다}}

\begin{document}
\maketitle

\begin{abstract}
직전 노트는 [달력제거] 판에서 「순서는 잡힌다」를 확정했다. 프레시 티처가 그 결론을
세 곳에서 무너뜨렸다 --- 「여는 방향」으로 지목한 격자 정체 천장은 \textbf{자기 행 암기}였고
(LOO 로 재면 $-0.457$), 200/200 을 통과했다는 관문은 \textbf{원리상 발화 불가}였으며,
크기 판정의 두 수가 \textbf{서로 다른 양}이었다.
이 노트는 티처가 지목한 판 --- \textbf{크기 물음이 살아 있는 유일한 판} --- 에서 다시 쟀다.
\claimx{\textbf{원판에서도 이겼다} --- $k=0/200$, $p=\mathbf{0.004975}$.
개선차 $\mathbf{653.1}$ 분.}
{진짜 $0.754520$ $\cdot$ 귀무 평균 $0.300972$ $\cdot$ SD $0.021173$ $\cdot$
귀무 최대 $0.355838$ 조차 진짜보다 한참 아래다.}
\claimx{🔴 \textbf{효과의 절반을 나르는 한 칸을 통째로 빼도 이겼다} ---
\texttt{b\_prv}$=-1$ 칸(첫 간격 행 $4.15\%$)을 없앤 판에서 $k=0/200$, $p=\mathbf{0.004975}$.}
{개선차 $\mathbf{347.1}$ 분. 🔴 그리고 그 칸을 빼는 것이 \textbf{어느 쪽에 불리한지 정했다} ---
판①$-$판② 차가 진짜($0.306646$)에서 귀무 평균($0.094172$)보다 크고 \textbf{200/200} 한 방향이다.
즉 \textbf{칸을 빼는 것은 진짜에게 불리(보수적)}하다. 「모른다」로 안 적었다.}
\claimx{🔴 \textbf{이번 관문은 원리상 걸릴 수 있었다} ---
귀무 기준 팔 MAE 의 가짓수가 $\mathbf{19/20}$ 이다.}
{직전 노트에서는 그 가짓수가 \textbf{1} 이라 상대차 집합이 $\{0.0\}$ 이었고,
「200/200 통과」는 통과가 아니라 \textbf{검정력 0} 이었다. 이 판은 기준 팔이 상수가 아니다.}
\claimx{🔴🔴 \textbf{그런데 크기는 또 측정 전에 졌다} ---
$Z=X/Y$ 가 판 ① $\mathbf{1.287}$, 판 ② $\mathbf{2.338}$ 로 \textbf{둘 다 1 을 넘는다}.}
{이번 $Y$ 는 \textbf{LOO 로 쟀다}(규율 4 개정판). 판 ① 의 진짜 개선은 그 천장의 \textbf{97.12\%} 다.
🔴 판정 눈금도 \textbf{개선} 하나로 통일했다 --- 효과 눈금은 병기다.}
\end{abstract}

\section{티처가 지목한 판에서 쟀다}

네 판의 $Z$ 는 원판 $1.155$ $\cdot$ 첫간격제외 $2.134$ $\cdot$ 달력제거 $3.045$ $\cdot$ 둘 다 $4.587$ 였다.
🔴 \textbf{원판이 1 에 가장 가깝고, 정답 다중집합이 같은 판인데 아무도 판정에 안 썼다.}

\begin{figure}[h]\centering
\includegraphics[width=\linewidth]{figs/fig1.pdf}
\caption{(a) 판 ① --- 원판 전량. (b) 판 ② --- \texttt{b\_prv}$=-1$ 칸을 통째로 뺐다.
둘 다 $k=0/200$. (c) 🔴 귀무 기준 팔 MAE 의 가짓수 --- 직전 노트는 \textbf{1}(발화 불가), 이번엔 \textbf{19}.}
\end{figure}

\begin{table}[h]\centering\small
\begin{tabular}{lrr}
\toprule
& 판 ① 원판 전량 & 판 ② 그 칸 제외 \\
\midrule
진짜 개선 & $0.754520$ & $0.447875$ \\
귀무 200 평균 & $0.300972$ & $0.206800$ \\
귀무 SD (ddof=1) & $0.021173$ & $0.014337$ \\
귀무 최대 & $0.355838$ & $0.246991$ \\
$k$ $\cdot$ $p$ & $\mathbf{0}$ $\cdot$ $\mathbf{0.004975}$ & $\mathbf{0}$ $\cdot$ $\mathbf{0.004975}$ \\
개선차 & \textbf{653.1 분} & \textbf{347.1 분} \\
$Y$ (LOO 천장) & $0.776895$ & $0.427689$ \\
$Z=X/Y$ & $\mathbf{1.287}$ & $\mathbf{2.338}$ \\
\bottomrule
\end{tabular}
\caption{🔴 판 ② 는 판 ① 과 \textbf{같은 귀무 세계}를 쓴다 --- \textbf{독립 확인 측정이 아니다}.
직전 노트가 그 오용으로 지적받았기에 이번엔 산출물이 그 사실을 필드로 들고 있다.}
\end{table}

\section{🔴 부호를 정했다 --- 「모른다」로 안 적었다}

「그 칸을 빼면 누가 손해인가」는 두 사이클 연속 \textbf{「못 심는 것」}에 들어가 있던 항목이다.
이번에 실측으로 정했다.

\begin{figure}[h]\centering
\includegraphics[width=\linewidth]{figs/fig2.pdf}
\caption{(a) 규율 4 개정판 --- $Y$ 를 \textbf{LOO} 로 쟀는데도 두 판 다 $Z>1$.
(b) 그 칸이 나르는 몫: 진짜가 귀무보다 큰 뽑기 \textbf{200/200} $\to$ \textbf{보수적}.
(c) 뽑기마다 찍은 검사와 순열이 부순 격자 비율.}
\end{figure}

판①$-$판② 의 차는 진짜에서 $0.306646$, 귀무 평균에서 $0.094172$ 이고
\textbf{200/200 뽑기에서 진짜가 크다}. 그리고 판 ② 에서 귀무의 남은 정답이 기후값에서
\textbf{더 멀다}(상대차 평균 $+0.036248$ $\cdot$ 양수인 뽑기 200/200).
그러므로 \textbf{칸을 빼는 것은 진짜에게 불리하고, 그럼에도 $k=0$ 이다.}

\section{🔴 그런데 크기는 또 측정 전에 졌다}

이번 $Y$ 는 \textbf{LOO}(자기 행 제외)로 쟀다 --- 직전 사이클에서 자기 행을 포함해 재는 바람에
「여는 방향」이 \textbf{암기}였던 것이 드러났기 때문이다.
그렇게 재도 $Z$ 가 판 ① $1.287$, 판 ② $2.338$ 로 \textbf{둘 다 1 을 넘는다}.
판 ① 의 진짜 개선은 이미 천장의 \textbf{97.12\%} 이고, 개선 눈금으로 \textbf{1.325 배} 모자란다.

⚠ 🔴 \textbf{판 ② 의 진짜 개선은 자기 $Y$ 의 104.72\% 다.}
LOO 천장은 \emph{추정}이지 상한이 아니므로 넘을 수 있다 --- \textbf{그 사실을 여기 적어 둔다}.
$Z$ 판정에는 영향이 없다($Z>1$ 은 두 판 다 유지).

\section{🔴 정정 --- 「대부분」이 아니라 「절반」이다}

이 논문의 초판 제목은 「효과의 \textbf{대부분}을 나르는 칸」이었다. 그것은 틀렸다.
그 87.5\% 는 \textbf{다른 판([달력제거])의 칸 분해}이고, \textbf{원판}에서는 이렇다 ---
행의 $4.15\%$ 인 그 칸이 \textbf{효과의 49.05\% $\cdot$ 개선의 43.10\%} 를 나른다
(칸 밖 $0.23106$ $+$ 칸 $0.22249$ $=$ $0.45355$ 로 \textbf{합이 전량과 비트 일치}).

⚠ 🔴 \textbf{두 정의가 갈린다는 것도 적어 둔다.} 「칸 분해」(행을 갈라 더하면 전량)는 $49.05\%$ 이고,
「판 차이」$(\text{판}①-\text{판}②)/\text{판}①$ 로 재면 $\mathbf{46.85\%}$ 다 --- 주 세션이 직접 셌다.
후자는 기준 팔의 분모까지 함께 바뀌므로 \textbf{같은 양이 아니다}.
어느 쪽이든 \textbf{「대부분」은 아니다.}

\section{남는 문장}

이 자료가 지지하는 문장은 이제 이것이다:
\textbf{「격자 안 간격 순서를 부수면 예측이 무너진다 --- 첫 간격 칸을 통째로 빼도 그렇다.」}
그리고 함께 남는 문장은 이것이다:
\textbf{「그 크기는 어느 판에서도 채택 크기에 못 미치고, 그것은 문턱이 아니라 조건 열의 성질이다.」}

\end{document}
